\documentclass[11pt]{article}
% ======================================================================
%  Preambule commun - Corriges Bac Serie A (Mathematiques)
%  Style sobre : noir et blanc, sans couleurs, sans filigrane.
% ======================================================================
\usepackage[utf8]{inputenc}
\usepackage[T1]{fontenc}
\usepackage{lmodern}
\usepackage{textcomp}
\usepackage{amsmath,amssymb}
\usepackage{enumitem}
\usepackage{array}
\usepackage[a4paper,margin=2.3cm]{geometry}
\usepackage{fancyhdr}
\usepackage{tikz}
\usetikzlibrary{arrows.meta,calc,patterns}

\DeclareUnicodeCharacter{00B0}{\ensuremath{^\circ}}
\tolerance=2000
\emergencystretch=2em
\frenchspacing

\pagestyle{fancy}
\fancyhf{}
\renewcommand{\headrulewidth}{0pt}
\fancyfoot[C]{\thepage}

\newcommand{\R}{\mathbb{R}}
\newcommand{\N}{\mathbb{N}}
\newcommand{\Z}{\mathbb{Z}}
\newcommand{\vi}{\vec{\imath}}
\newcommand{\vj}{\vec{\jmath}}

\newcommand{\exo}[2]{\bigskip\noindent{\large\bfseries Exercice #1}\hfill\textit{#2}\par\nopagebreak\smallskip}
\newcommand{\partie}[1]{\medskip\noindent{\bfseries Partie #1}\par\nopagebreak\smallskip}
\newcommand{\entetesujet}[1]{\begin{center}{\Large\bfseries #1}\end{center}\vspace{0.3em}\hrule\vspace{1.1em}}

\newlist{questions}{enumerate}{1}
\setlist[questions]{label=\textbf{\arabic*.},leftmargin=1.8em,itemsep=3pt,topsep=3pt}
\newlist{sousq}{enumerate}{1}
\setlist[sousq]{label=\textbf{\alph*.},leftmargin=1.8em,itemsep=2pt,topsep=2pt}


\begin{document}

\entetecorrige{2021}{A}{Mathématiques}

\exo{1}{}
\partie{A}
\begin{questions}
  \item $a = \dfrac{9}{10}$ et $b = \dfrac{18}{25}$ sont déjà irréductibles (seul diviseur commun : 1). $c = \dfrac{4}{30} = \dfrac{2\times 2}{2\times 15} = \dfrac{2}{15}$.
  \item $\mathrm{PPCM}(10, 25, 15)$ : en décomposant, $10 = 2\times 5$, $25 = 5^2$, $15 = 3\times 5$. Le PPCM est le produit des facteurs premiers, chacun avec son plus grand exposant : $2\times 5^2\times 3 = 150$.
  \item \begin{sousq}
    \item Réduction au dénominateur 150 : $\dfrac{9}{10} = \dfrac{135}{150}$ ($\times 15$) ; $\dfrac{18}{25} = \dfrac{108}{150}$ ($\times 6$) ; $\dfrac{2}{15} = \dfrac{20}{150}$ ($\times 10$).
    \item $S = \dfrac{9}{10} - \dfrac{18}{25} + \dfrac{4}{30} = \dfrac{135}{150} - \dfrac{108}{150} + \dfrac{20}{150} = \dfrac{135 - 108 + 20}{150} = \dfrac{47}{150}$.
  \end{sousq}
\end{questions}

\partie{B}
\begin{questions}
  \item \begin{sousq}
    \item $P(x) = (2x+1)(x^2 + x - 6) = 2x^3 + 2x^2 - 12x + x^2 + x - 6 = 2x^3 + 3x^2 - 11x - 6$.
    \item $x^2 + x - 6 = 0$ : $\Delta = 1 - 4(-6) = 25$, d'où $x_1 = \dfrac{-1 - 5}{2} = -3$ et $x_2 = \dfrac{-1 + 5}{2} = 2$.
    \item $P(x) = 0 \Leftrightarrow 2x + 1 = 0$ ou $x^2 + x - 6 = 0$, soit $x = -\dfrac{1}{2}$, $x = -3$ ou $x = 2$. $S = \left\{-3\,;-\dfrac{1}{2}\,;2\right\}$.
  \end{sousq}
  \item En posant $x = \ln t$, $(E) : 2(\ln t)^3 + 3(\ln t)^2 - 11\ln t - 6 = 0$ devient $2x^3 + 3x^2 - 11x - 6 = 0$, de solutions $\{-3\,;-\tfrac{1}{2}\,;2\}$. D'où $\ln t_1 = -3$, $\ln t_2 = -\tfrac{1}{2}$, $\ln t_3 = 2$, soit $t \in \left\{e^{-3}\,;e^{-\frac{1}{2}}\,;e^{2}\right\}$.
\end{questions}

\exo{2}{}
\begin{questions}
  \item \begin{sousq}
    \item $\R^* = ]-\infty;0[\,\cup\,]0;+\infty[$.
    \item $\dfrac{x-1}{x} = \dfrac{x}{x} + \dfrac{-1}{x} = 1 + \dfrac{-1}{x}$, donc $a = -1$.
    \item $\lim\limits_{x\to-\infty} g(x) = \lim\limits_{x\to-\infty}\dfrac{x}{x} = 1$ : asymptote horizontale $y = 1$ en $-\infty$. $\lim\limits_{x\to 0^-} g = \lim\limits_{x\to 0^-}\dfrac{-1}{x} = +\infty$ et $\lim\limits_{x\to 0^+} g = -\infty$ : asymptote verticale $x = 0$.
  \end{sousq}
  \item \begin{sousq}
    \item $\forall x \in \R^*$, $g'(x) = \dfrac{x - (x-1)}{x^2} = \dfrac{1}{x^2}$.
    \item $g'(x) = \dfrac{1}{x^2} > 0$ : $g$ est strictement croissante sur $\R^*$.
    \item \begin{center}\renewcommand{\arraystretch}{1.4}
    \begin{tabular}{|c|ccccc|}
    \hline $x$ & $-\infty$ & & $0$ & & $+\infty$ \\\hline
    $g'(x)$ & & $+$ & $\|$ & $+$ & \\\hline
    $g$ & $1$ & $\nearrow$ & $+\infty\ \|\ -\infty$ & $\nearrow$ & $1$ \\\hline
    \end{tabular}
    \end{center}
  \end{sousq}
  \item \begin{sousq}
    \item $g(-1) = \dfrac{-1-1}{-1} = 2$ : $A(-1\,;2)$. $g(1) = 0$ : $B(1\,;0)$. $g(2) = \dfrac{2-1}{2} = \dfrac{1}{2}$ : $C\left(2\,;\dfrac{1}{2}\right)$.
    \item
    \begin{center}
    \begin{tikzpicture}[>=Stealth,scale=0.5,every node/.style={font=\scriptsize}]
      \draw[->] (-7.3,0)--(7.3,0) node[right]{$x$};
      \draw[->] (0,-5.3)--(0,5.3) node[above]{$y$};
      \foreach \x in {-7,-6,-5,-4,-3,-2,-1,1,2,3,4,5,6,7}{\draw (\x,0.1)--(\x,-0.1) node[below]{\x};}
      \foreach \y in {-5,-4,-3,-2,-1,1,2,3,4,5}{\draw (0.1,\y)--(-0.1,\y) node[left]{\y};}
      \draw[densely dashed] (-7.3,1)--(7.3,1);
      \draw[thick] plot[domain=0.17:7,samples=120] (\x,{max(-5.2,1-1/\x)});
      \draw[thick] plot[domain=-7:-0.17,samples=120] (\x,{min(5.2,1-1/\x)});
      \fill (-1,2) circle (3pt) node[above left]{$A$};
      \fill (1,0) circle (3pt) node[below right]{$B$};
      \fill (2,0.5) circle (3pt) node[above right]{$C$};
      \node at (4.5,1.5){$(\mathcal{C}_g)$};
    \end{tikzpicture}
    \end{center}
  \end{sousq}
\end{questions}

\exo{3}{}
\begin{questions}
  \item $\overline{X} = \dfrac{1+2+3+4+5+6}{6} = 3{,}5$ ; $\overline{Y} = \dfrac{1{,}2 + 1{,}8 + 2{,}3 + 2{,}5 + 3 + 3{,}5}{6} = 2{,}38$. Point moyen $G(3{,}5\,;2{,}38)$.
  \item \begin{sousq}
    \item Sous-tableaux :
    \begin{center}\renewcommand{\arraystretch}{1.3}
    $T_1$ : \begin{tabular}{|c|c|c|c|}\hline $X$ (mois) & 1 & 2 & 3 \\\hline $Y$ (mg) & 1{,}2 & 1{,}8 & 2{,}3 \\\hline\end{tabular}\qquad
    $T_2$ : \begin{tabular}{|c|c|c|c|}\hline $X$ (mois) & 4 & 5 & 6 \\\hline $Y$ (mg) & 2{,}5 & 3 & 3{,}5 \\\hline\end{tabular}
    \end{center}
    \item $\overline{X}_1 = \dfrac{1+2+3}{3} = 2$, $\overline{Y}_1 = \dfrac{1{,}2 + 1{,}8 + 2{,}3}{3} = 1{,}77$ : $G_1(2\,;1{,}77)$. $\overline{X}_2 = \dfrac{4+5+6}{3} = 5$, $\overline{Y}_2 = \dfrac{2{,}5 + 3 + 3{,}5}{3} = 3$ : $G_2(5\,;3)$.
    \item Droite de Mayer (par $G_1$, $G_2$) : $a = \dfrac{3 - 1{,}77}{5 - 2} = 0{,}41$. Comme $G_2 \in$ la droite : $3 = 0{,}41\times 5 + b$, d'où $b = 0{,}95$. Ainsi $y = 0{,}41x + 0{,}95$.
  \end{sousq}
  \item En multipliant par 3 : $3y = 1{,}23x + 2{,}85$, soit $1{,}23x - 3y + 2{,}85 = 0$, quasi identique à $1{,}24x - 3y + 2{,}8 = 0$. Pour $x = 7$ : $1{,}24\times 7 - 3y + 2{,}8 = 0 \Rightarrow 3y = 11{,}48 \Rightarrow y \approx 3{,}83$. Au 7\textsuperscript{e} mois, le poids de la larve est d'environ 3,83 mg.
\end{questions}

\end{document}
